\section{Conclusion}\label{sec:conclusion}

When a husband loses his job, his wife is ten times as likely to lose hers in the same month. The excess is carried by
shared workplaces, not by the local business cycle, and it is the same in Black and White couples. It is small in the
household's average income loss but large in the month of the shock and among couples who share a workplace, which is
where insurance is valued most. Wives who were not working search at once and find jobs slowly, so in the month in which
the household's income falls most its second earner is more likely to lose her own job than to have added work. Average
added-worker estimates net these two groups of wives. They understate what wives do and overstate what families can
absorb.

The fact ranks the safety net. A transfer that requires work, such as the Earned Income Tax Credit, pays nothing in the
months in which both spouses are out of work, and over the year it rises by \$196 on average for couples who lose
both jobs, while food stamps pay in every state. The model of
Section~\ref{sec:model} is built to measure how much this matters for consumption and for children, and how the
ranking of work-conditioned and unconditional transfers changes when the second earner's job is not an independent
source of insurance.
